% ECONOMICS_PROBLEM: {"id": "C23-429", "title": "Sparse synthetic controls", "jel_code": "C23", "source_id": "OP-0425"}
% !TeX program = xelatex
\documentclass[11pt,a4paper]{article}
\usepackage[margin=23mm]{geometry}
\usepackage{fontspec}
\setmainfont{Latin Modern Roman}
\usepackage{amsmath,amssymb,mathtools}
\usepackage{xcolor,enumitem,booktabs,longtable,array,xurl,hyperref}
\definecolor{muted}{HTML}{53636C}
\hypersetup{colorlinks=true,urlcolor=blue,linkcolor=blue}
\setlength{\parindent}{0pt}
\setlength{\parskip}{5pt}
\newcommand{\R}{\mathbb R}
\newcommand{\E}{\mathbb E}
\newcommand{\N}{\mathbb N}
\newcommand{\ind}{\mathbf 1}
\newcommand{\Var}{\operatorname{Var}}
\newcommand{\Cov}{\operatorname{Cov}}
\newcommand{\supp}{\operatorname{supp}}
\DeclareMathOperator*{\argmin}{arg\,min}
\DeclareMathOperator*{\argmax}{arg\,max}
\newcommand{\entrymeta}[3]{{\sffamily\small\color{muted}\textbf{\texttt{#1}}\quad|\quad #2\par #3\par}}
\newcommand{\Problemref}[1]{\texttt{#1}}
\newcommand{\JELref}[2]{\texttt{#2} (JEL \texttt{#1})}
\begin{document}
\section*{Sparse synthetic controls}
\textbf{Problem ID:} \texttt{C23-429}\quad \textbf{JEL:} \texttt{C23}\par
% BEGIN ECONOMICS STATEMENT
\label{problem:OP-0425}
{\sffamily\small\color{muted}Original subject group: Econometrics, causal inference and measurement\par}
\label{definitions:problem:30007707}
\textbf{Audit classification:} Published research agenda. The source research agenda is verified at p.423. Modern inference results exist; this finite-donor robust specification is editorial, and narrow coverage cannot follow solely from pretreatment fit and a rank bound.
\subsection*{Definitions and precise research target}
Consider one treated region \(j=0\), \(J\) donor regions, \(T_0\) pretreatment periods and a fixed post-treatment horizon \(H\). Untreated outcomes satisfy \(Y_{jt}(0)=\lambda_j^\top f_t+\eta_{jt}+\xi_{jt}\), where \(\lambda_j,f_t\in\mathbb R^r\), the known rank bound is \(r\leq r_{\max}\), and their Euclidean norms are bounded by stated constants \(\Lambda,F\). Approximation errors obey \(|\eta_{jt}|\leq a_{jt}\); noises \(\xi_{jt}\) are mean-zero and independent across regions and periods, with a stated sub-Gaussian scale bound \(\sigma\). Donor weights \(w_j\geq0\) sum to one and are selected using pretreatment data. Post-treatment donor spillovers satisfy \(|s_{jt}|\leq b_{jt}\), for prespecified nonnegative bounds \(a_{jt},b_{jt}\). The target is \(\tau=H^{-1}\sum_{t=T_0+1}^{T_0+H}\{Y_{0t}(1)-Y_{0t}(0)\}\), the treated region's average effect. For fixed small \(J\), construct intervals with coverage at least \(1-\alpha\), for specified \(\alpha\in(0,1)\), uniformly over this model class, including weight selection. Determine when pretreatment fit and validation periods make intervals informative, and how removing contaminated donors trades spillover protection against worse fit. No inference guarantee is asserted without the declared factor, approximation and noise restrictions. The source's broad agenda supports inference and sensitivity concerns; this finite-donor specification is an adaptation whose contemporary unresolved status needs separate validation.

\textbf{Definitions, assumptions and mathematical qualifications.} The treated region is observed under treatment post-policy; donors satisfy \(Y_{jt}=Y_{jt}(0)+s_{jt}\) after treatment, and all regions are untreated before it. Weights are measurable functions of pretreatment observations only. Conditional on deterministic factor primitives, \(E[e^{u\xi_{jt}}]\leq e^{u^2\sigma^2/2}\) for every real \(u\); clarify any dependence of treatment outcomes on untreated noise. Uniform coverage is \(\inf_{M\in\mathcal M}\Pr_M(\tau_M\in C)\geq1-\alpha\), including weight randomness and counterfactual prediction noise. Pretreatment fit alone does not identify post-treatment factors: a factor direction absent before treatment can emerge afterward. Informative narrow intervals require a stated extrapolation, span or factor-excitation condition, not only bounded rank.
\subsection*{Variants and assumption changes}
\begin{enumerate}
\item \textbf{Editorial, exact span.} Assume \(\lambda_0=\sum_jw_j\lambda_j\) for admissible nonnegative weights and no donor spillovers; account for estimated weights, approximation error and post-treatment noise.
\item \textbf{Editorial, bounded contamination.} Permit \(|s_{jt}|\leq b_{jt}\) and loading mismatch \(\|\lambda_0-\sum_jw_j\lambda_j\|\leq d(w)\). Compare intervals using the worst-case factor bias \(F d(w)\), approximation bounds and weighted spillover costs.
\end{enumerate}
\subsection*{References and source audit}
\textbf{Checked source.} Journal of Economic Literature 59(2) (2021), pp. 391--425, Section 9 p. 423; Section 6 pp. 410--411. The conclusion identifies inference and sensitivity as research areas. Full primary text retrieved. \url{https://economics.mit.edu/sites/default/files/publications/jel.20191450.pdf}.
\begin{enumerate}\raggedright
\item\hypertarget{definitions:source:30007707:1}{} Alberto Abadie. 2021. \emph{Using Synthetic Controls: Feasibility, Data Requirements, and Methodological Aspects}. Journal of Economic Literature 59(2) (2021), pp. 391--425, Section 9 p. 423; Section 6 pp. 410--411.
\url{https://economics.mit.edu/sites/default/files/publications/jel.20191450.pdf}.
DOI: \url{https://doi.org/10.1257/jel.20191450}.
\end{enumerate}

\subsection*{Common conventions: Source mathematical conventions}

These conventions apply to each entry unless explicitly replaced.

\textbf{Models and identification.} A primitive model \(m\) specifies all probability laws, feasible choices, information, equilibrium and measurement rules used by its target. The admissible class \(\mathcal M\) is fixed independently of outcomes. If \(P_O(m)\) is its observable probability law and \(\tau(m)\) the target, its identified set at observable law \(P\) is
\[
 \mathcal I_\tau(P)=\{\tau(m):m\in\mathcal M,\ P_O(m)=P\}.
\]
Point identification means that this set is a singleton. An empty set signals incompatibility of data and assumptions. Coordinate bounds need not describe the joint set. Bounds are sharp only if every retained target is attainable or arbitrarily approachable by compatible models. Finite bounds require explicit restrictions on unobserved quantities; unrestricted confounding, hidden wealth, extrapolation or arbitrary equilibrium selection can yield uninformative sets.

\textbf{Causal interventions.} A potential outcome \(Y_i(p)\) is the outcome under a complete policy regime \(p\), including timing, financing, information and market spillovers. The target population and units are fixed before treatment. With interference, use the complete assignment vector or a justified exposure mapping; individual no-interference notation alone cannot identify an economy-wide intervention. A difference of mean potential outcomes does not require the cross-world joint distribution. Conditioning on post-treatment migration, survival, enrollment or employment changes the estimand and needs separate justification.

\textbf{Statistical statements.} An estimator, confidence set, forecast or test specifies a sampling law, model class and loss/coverage criterion. Uniform \((1-\alpha)\)-coverage means \(\inf_{m\in\mathcal M}\Pr_m\{\tau(m)\in C_n\}\geq1-\alpha\), or an explicitly asymptotic version. The whole parameter space trivially has coverage, so informativeness requires a stated width, risk or power objective. A forecasting improvement is relative to a named benchmark, information set and evaluation population.

\textbf{Equilibrium and optimization.} An equilibrium comprises feasible plans, optimality, beliefs where needed, budget constraints and all market/resource clearing. The primitives or a stated benchmark define each correspondence \(\mathcal E(p;m)\). Nonemptiness is assumed or proved, never inferred just from compact policy sets. A selected-equilibrium target uses a declared selection rule; a robust target quantifies over all permitted equilibria. Use suprema or infima unless attainment is established. An \(\varepsilon\)-optimal policy is feasible and within \(\varepsilon\) of the finite optimum value. Report an empty feasible set separately.

\textbf{Welfare and accounting.} Normative utility, interpersonal normalization, social weights, numeraire, discounting and the use of public receipts are primitives. Behavioral utility need not coincide with the welfare criterion. Household welfare includes ownership income and rents once; adding profits again to compensating variation generally double-counts them. A monetary equivalent is defined by an explicit transfer and price convention, with monotonicity and range conditions. Average effects, mechanism contributions, transfers and welfare losses are different targets. Infinite sums require convergence and dynamic feasible plans require terminal or no-Ponzi conditions.

\textbf{Source versions.} A working-paper date, revision date and journal publication year are distinguished. A DOI for a journal version is not automatically the DOI of an earlier manuscript. Locators refer to the stated version and printed pages where specified. A metadata-only check does not verify a claim about a full-text conclusion. Every entry's source subsection narrows the provenance claim accordingly.
% END ECONOMICS STATEMENT
\end{document}
